% =================================================
% CHAPITRE 1 : INTRODUCTION
% =================================================

\chapter{Introduction}
\label{chap-introduction}

\section{Contexte}

MiniNumpy est une bibliothèque de calcul numérique
écrite de zéro en C++20, sans dépendance extérieure.
Elle reprend, à petite échelle, ce que fait NumPy
en Python~\cite{numpy} : manipuler des vecteurs et
des matrices de réels et leur appliquer des
opérations.

Le projet a un but final précis : réimplémenter en
C++, avec MiniNumpy à la place de NumPy, le réseau
de neurones du livre \emph{Neural Networks from
Scratch in Python}~\cite{nnfs}, puis l'entraîner sur
le jeu de données Fashion-MNIST (des images de
vêtements à classer en dix catégories).

C'est ce but qui fixe le contenu de ce carnet :
chaque notion présentée ici sert, à un moment ou
à un autre, à faire fonctionner ce réseau.

\section{Rôle de ce carnet}

Écrire une bibliothèque numérique oblige à
comprendre exactement ce que l'on calcule : une
formule mal comprise donne un code qui compile mais
renvoie de faux résultats. Ce carnet fait donc le
lien entre la théorie et le code.

Chaque notion y est présentée de la même façon :
\begin{enumerate}
    \item une \textbf{définition} précise ;
    \item un \textbf{exemple calculé à la main}, sur
          de petites valeurs, qui sert aussi de cas
          de test ;
    \item son \textbf{utilité dans le projet} :
          à quoi elle sert dans le réseau de
          neurones ;
    \item son \textbf{implémentation} : la classe
          ou la fonction de MiniNumpy qui la
          réalise.
\end{enumerate}

\section{Plan}

Le carnet suit l'ordre dans lequel la bibliothèque
est construite. Les chapitres déjà rédigés :
\begin{enumerate}
    \item \textbf{Vecteurs} (\cref{chap-vecteurs}).
          Les opérations de base, le produit
          scalaire et la norme. Un neurone calcule
          un produit scalaire entre ses entrées et
          ses poids.
    \item \textbf{Matrices} (\cref{chap-matrices}).
          Stockage, accès aux coefficients,
          transposée et produit matriciel. Les poids
          d'une couche entière forment une matrice,
          et un seul produit matriciel calcule la
          sortie de tous les neurones de la couche,
          pour tout un lot d'images.
    \item \textbf{Systèmes linéaires}
          (\cref{chap-systemes}). Pivot de Gauss,
          décomposition LU, déterminant et inverse.
          Ils servent aux moindres carrés et à
          l'optimisation.
\end{enumerate}

Les chapitres à venir suivent les étapes du projet :
moindres carrés, aléatoire et Monte Carlo,
optimisation, valeurs propres, différentiation
automatique. Viennent enfin les notions propres au
réseau de neurones :
\begin{enumerate}[resume]
    \item \textbf{Broadcasting.} La règle qui permet
          d'additionner un vecteur à chaque ligne
          d'une matrice, par exemple pour ajouter
          les biais d'une couche.
    \item \textbf{Fonctions d'activation.} ReLU et
          ses voisines, qui rendent le réseau non
          linéaire.
    \item \textbf{Softmax.} Elle transforme les
          sorties de la dernière couche en
          probabilités pour les dix classes.
    \item \textbf{Entropie croisée.} La fonction de
          perte, qui mesure l'écart entre la
          prédiction du réseau et la bonne réponse.
    \item \textbf{Rétropropagation.} Dérivées
          partielles, gradient et dérivation en
          chaîne : comment chaque poids influe sur
          l'erreur.
    \item \textbf{Descente de gradient.} Elle
          corrige les poids pas à pas pour faire
          baisser l'erreur ; c'est l'entraînement
          proprement dit.
\end{enumerate}

\section{Conventions de notation}

\begin{itemize}
    \item Les scalaires sont notés en minuscules
          italiques : $a$, $\lambda$, $x_i$.
    \item Les vecteurs sont notés en minuscules
          grasses : $\mathbf{x} \in \R^n$. Ce sont
          des vecteurs colonnes.
    \item Les matrices sont notées en majuscules :
          $W \in \R^{m \times n}$ a $m$ lignes et
          $n$ colonnes, et $w_{ij}$ est le
          coefficient de la ligne $i$ et de la
          colonne $j$.
    \item La transposée de $W$ est notée $W^\top$.
    \item En mathématiques, les indices commencent
          à $1$ : $\mathbf{x} = (x_1, \dots, x_n)$.
          En C++, ils commencent à $0$ : la
          composante $x_i$ correspond donc à
          \cpp|v[i-1]| dans le code.
\end{itemize}


